% Options for packages loaded elsewhere
\PassOptionsToPackage{unicode}{hyperref}
\PassOptionsToPackage{hyphens}{url}
\PassOptionsToPackage{dvipsnames,svgnames,x11names}{xcolor}
%
\documentclass[
  11pt,
]{article}
\usepackage{amsmath,amssymb}
\usepackage{iftex}
\ifPDFTeX
  \usepackage[T1]{fontenc}
  \usepackage[utf8]{inputenc}
  \usepackage{textcomp} % provide euro and other symbols
\else % if luatex or xetex
  \usepackage{unicode-math} % this also loads fontspec
  \defaultfontfeatures{Scale=MatchLowercase}
  \defaultfontfeatures[\rmfamily]{Ligatures=TeX,Scale=1}
\fi
\usepackage[]{lmodern}
\ifPDFTeX\else
  % xetex/luatex font selection
\fi
% Use upquote if available, for straight quotes in verbatim environments
\IfFileExists{upquote.sty}{\usepackage{upquote}}{}
\IfFileExists{microtype.sty}{% use microtype if available
  \usepackage[]{microtype}
  \UseMicrotypeSet[protrusion]{basicmath} % disable protrusion for tt fonts
}{}
\makeatletter
\@ifundefined{KOMAClassName}{% if non-KOMA class
  \IfFileExists{parskip.sty}{%
    \usepackage{parskip}
  }{% else
    \setlength{\parindent}{0pt}
    \setlength{\parskip}{6pt plus 2pt minus 1pt}}
}{% if KOMA class
  \KOMAoptions{parskip=half}}
\makeatother
\usepackage{xcolor}
\usepackage[margin=2.2cm]{geometry}
\usepackage{listings}
\newcommand{\passthrough}[1]{#1}
\lstset{defaultdialect=[5.3]Lua}
\lstset{defaultdialect=[x86masm]Assembler}
\usepackage{longtable,booktabs,array}
\usepackage{calc} % for calculating minipage widths
% Correct order of tables after \paragraph or \subparagraph
\usepackage{etoolbox}
\makeatletter
\patchcmd\longtable{\par}{\if@noskipsec\mbox{}\fi\par}{}{}
\makeatother
% Allow footnotes in longtable head/foot
\IfFileExists{footnotehyper.sty}{\usepackage{footnotehyper}}{\usepackage{footnote}}
\makesavenoteenv{longtable}
\setlength{\emergencystretch}{3em} % prevent overfull lines
\providecommand{\tightlist}{%
  \setlength{\itemsep}{0pt}\setlength{\parskip}{0pt}}
\setcounter{secnumdepth}{-\maxdimen} % remove section numbering
\usepackage{xcolor}
\definecolor{codebg}{RGB}{246,248,250}
\lstset{basicstyle=\ttfamily\small,breaklines=true,breakatwhitespace=false,
  backgroundcolor=\color{codebg},frame=single,rulecolor=\color{gray!40},
  columns=fullflexible,keepspaces=true,showstringspaces=false,upquote=true,
  postbreak=\mbox{\textcolor{gray}{$\hookrightarrow$}\space},
  literate={ğ}{{\u{g}}}1 {ç}{{\c{c}}}1 {ı}{{\i}}1 {ş}{{\c{s}}}1 {ö}{{\"o}}1 {ü}{{\"u}}1}
\usepackage{etoolbox}
\AtBeginEnvironment{longtable}{\footnotesize}
\setlength{\emergencystretch}{3em}
\sloppy
\ifLuaTeX
  \usepackage{selnolig}  % disable illegal ligatures
\fi
\usepackage{bookmark}
\IfFileExists{xurl.sty}{\usepackage{xurl}}{} % add URL line breaks if available
\urlstyle{same}
\hypersetup{
  colorlinks=true,
  linkcolor={Maroon},
  filecolor={Maroon},
  citecolor={Blue},
  urlcolor={Blue},
  pdfcreator={LaTeX via pandoc}}

\title{Destructuring Field Guide}
\author{Kamil Postaağası, 2023501030}
\date{}

\begin{document}
\maketitle

{
\hypersetup{linkcolor=}
\setcounter{tocdepth}{3}
\tableofcontents
}
\section{Destructuring Field Guide}\label{destructuring-field-guide}

Name: Kamil Postaağası\\
Student number: 2023501030

Bonuses claimed: \textbf{Bonus 1: LaTeX}, \textbf{Bonus 2: Depth},
\textbf{Bonus 3: Real code}, \textbf{Bonus 4: Git repository} (see the
sections at the end).

\subsection{Run}\label{run}

\begin{lstlisting}
npm install
npm run all      # runs the JS and TS examples (output is in output.txt)
npm run check    # type-checks in strict mode
\end{lstlisting}

My examples use my own data: BIST stock prices, my portfolio positions,
my courses this term. Every output below is copied from
\passthrough{\lstinline!output.txt!}
(\passthrough{\lstinline!npm run all > output.txt!}).

\subsection{Field guide}\label{field-guide}

\subsubsection{\texorpdfstring{Form 1: values by position,
\texttt{const\ {[}x,\ y{]}\ =\ f()}}{Form 1: values by position, const {[}x, y{]} = f()}}\label{form-1-values-by-position-const-x-y-f}

\textbf{Syntax:}
\passthrough{\lstinline!const [lowest, highest] = minMax(closes)!}

\textbf{Example} (\passthrough{\lstinline!js/arrays.mjs!}):

\begin{lstlisting}
const closes = [41.2, 39.8, 43.5, 40.1]
const [lowest, highest] = minMax(closes)
console.log('form 1:', lowest, highest)
\end{lstlisting}

\textbf{Output:} \passthrough{\lstinline!form 1: 39.8 43.5!}

\textbf{Use it when:} a function gives back a few values in a fixed
order (like \passthrough{\lstinline!useState!} does).

\textbf{Watch out:} only the position counts, not the name. If I write
\passthrough{\lstinline![highest, lowest]!} the names are swapped but
the values are not. Without destructuring it is
\passthrough{\lstinline!const pair = minMax(closes); const lowest = pair[0]; const highest = pair[1]!}.

\subsubsection{\texorpdfstring{Form 2: skip and collect,
\texttt{const\ {[},\ second,\ ...others{]}\ =\ list}}{Form 2: skip and collect, const {[}, second, ...others{]} = list}}\label{form-2-skip-and-collect-const-second-...others-list}

\textbf{Syntax:}
\passthrough{\lstinline!const [, secondCourse, ...otherCourses] = courses!}

\textbf{Example} (\passthrough{\lstinline!js/arrays.mjs!}):

\begin{lstlisting}
const courses = ['MIS 233', 'INTT 321', 'MIS 251', 'INTT 375']
const [, secondCourse, ...otherCourses] = courses
\end{lstlisting}

\textbf{Output:}
\passthrough{\lstinline!form 2: INTT 321 [ 'MIS 251', 'INTT 375' ]!}

\textbf{Use it when:} I want one item from a known place and ``the
rest'' as a list.

\textbf{Watch out:} the rest element must be last
(\passthrough{\lstinline![...a, b]!} is a SyntaxError).
\passthrough{\lstinline!...!} here is \emph{rest} (it collects into a
new array); in \passthrough{\lstinline![...list]!} on the right side it
is \emph{spread} (it unpacks). Same three dots, opposite jobs.

\subsubsection{\texorpdfstring{Form 3: swap,
\texttt{{[}a,\ b{]}\ =\ {[}b,\ a{]}}}{Form 3: swap, {[}a, b{]} = {[}b, a{]}}}\label{form-3-swap-a-b-b-a}

\textbf{Syntax:} \passthrough{\lstinline![buy, sell] = [sell, buy]!}

\textbf{Example} (\passthrough{\lstinline!js/arrays.mjs!}):

\begin{lstlisting}
let buy = 'AKBNK'
let sell = 'GARAN'
;[buy, sell] = [sell, buy]
\end{lstlisting}

\textbf{Output:} \passthrough{\lstinline!form 3: GARAN AKBNK!}

\textbf{Use it when:} two variables need to trade values, no temp
variable needed.

\textbf{Watch out:} the variables must be \passthrough{\lstinline!let!}.
If the line before has no semicolon, JavaScript glues
\passthrough{\lstinline![!} to it
(\passthrough{\lstinline!sell = 'GARAN'[buy, sell] = ...!}), so I start
the line with \passthrough{\lstinline!;!}.

\subsubsection{\texorpdfstring{Form 4: array defaults,
\texttt{const\ {[}a\ =\ 1,\ b\ =\ 2{]}\ =\ {[}{]}}}{Form 4: array defaults, const {[}a = 1, b = 2{]} = {[}{]}}}\label{form-4-array-defaults-const-a-1-b-2}

\textbf{Syntax:}
\passthrough{\lstinline!const [currency = 'TRY', amount = 0] = []!}

\textbf{Example} (\passthrough{\lstinline!js/arrays.mjs!}):

\begin{lstlisting}
const [currency = 'TRY', amount = 0] = []
const [cur2 = 'TRY', amt2 = 0] = ['USD', undefined]
\end{lstlisting}

\textbf{Output:} \passthrough{\lstinline!form 4: TRY 0!} and
\passthrough{\lstinline!form 4b: USD 0!}

\textbf{Use it when:} the array can be shorter than I expect.

\textbf{Watch out:} the default is used only for
\passthrough{\lstinline!undefined!} (missing or written as undefined),
not for \passthrough{\lstinline!null!}, \passthrough{\lstinline!0!} or
\passthrough{\lstinline!''!}.

\subsubsection{\texorpdfstring{Form 5: properties by name,
\texttt{let\ \{\ a,\ b\ \}\ =\ t}}{Form 5: properties by name, let \{ a, b \} = t}}\label{form-5-properties-by-name-let-a-b-t}

\textbf{Syntax:}
\passthrough{\lstinline!let \{ ticker, shares \} = position!}

\textbf{Example} (\passthrough{\lstinline!js/objects.mjs!}):

\begin{lstlisting}
let { ticker, shares } = position
shares = shares + 10
console.log('form 5b:', shares, position.shares)
\end{lstlisting}

\textbf{Output:} \passthrough{\lstinline!form 5: ASELS 120!} and
\passthrough{\lstinline!form 5b: 130 120!}

\textbf{Use it when:} I need a few fields of an object and want short
names.

\textbf{Watch out:} the variable name must match the property name.
Changing \passthrough{\lstinline!shares!} changes only my variable, the
object still has 120 (for a primitive value; see pitfall 4 for objects).

\subsubsection{\texorpdfstring{Form 6: rename,
\texttt{const\ \{\ a:\ alpha\ \}\ =\ t}}{Form 6: rename, const \{ a: alpha \} = t}}\label{form-6-rename-const-a-alpha-t}

\textbf{Syntax:}
\passthrough{\lstinline!const \{ avgCost: costPerShare \} = position!}

\textbf{Example} (\passthrough{\lstinline!js/objects.mjs!}):

\begin{lstlisting}
const { avgCost: costPerShare } = position
console.log('form 6:', costPerShare)
\end{lstlisting}

\textbf{Output:} \passthrough{\lstinline!form 6: 58.4!}

\textbf{Use it when:} the property name is unclear here or already used
by another variable.

\textbf{Watch out:} after this there is no
\passthrough{\lstinline!avgCost!} variable, only
\passthrough{\lstinline!costPerShare!}. Left of the colon is where it
comes from, right of the colon is the new name. In TypeScript this looks
like a type annotation but it is not (topic T1).

\subsubsection{\texorpdfstring{Form 7: object defaults,
\texttt{const\ \{\ a\ =\ 1\ \}\ =\ t}}{Form 7: object defaults, const \{ a = 1 \} = t}}\label{form-7-object-defaults-const-a-1-t}

\textbf{Syntax:}
\passthrough{\lstinline!const \{ stopLoss = 50, shares: qty = 0 \} = position!}

\textbf{Example} (\passthrough{\lstinline!js/objects.mjs!}):

\begin{lstlisting}
const { stopLoss = 50, shares: qty = 0 } = position
\end{lstlisting}

\textbf{Output:} \passthrough{\lstinline!form 7: 50 120!}

\textbf{Use it when:} an optional setting should have a sensible value.

\textbf{Watch out:} \passthrough{\lstinline!shares!} exists (120), so
its default 0 is ignored. Rename and default can be combined:
\passthrough{\lstinline!shares: qty = 0!}. Again only
\passthrough{\lstinline!undefined!} triggers the default (pitfall 2).

\subsubsection{\texorpdfstring{Form 8: rest in an object,
\texttt{const\ \{\ a,\ ...others\ \}\ =\ t}}{Form 8: rest in an object, const \{ a, ...others \} = t}}\label{form-8-rest-in-an-object-const-a-...others-t}

\textbf{Syntax:}
\passthrough{\lstinline!const \{ owner, trades, ...summary \} = position!}

\textbf{Example} (\passthrough{\lstinline!js/objects.mjs!}):

\begin{lstlisting}
const { owner, trades, ...summary } = position
console.log('form 8:', summary)
\end{lstlisting}

\textbf{Output:}
\passthrough{\lstinline!form 8: \{ ticker: 'ASELS', shares: 120, avgCost: 58.4 \}!}

\textbf{Use it when:} I want to remove a few fields and keep everything
else, for example taking \passthrough{\lstinline!children!} out of props
and passing the rest on (P4).

\textbf{Watch out:} \passthrough{\lstinline!summary!} is a new object
but only a shallow copy (pitfall 4). Rest must be last.

\subsubsection{\texorpdfstring{Form 9: nested,
\texttt{const\ \{\ p:\ \{\ x\ \},\ list:\ {[}first{]}\ \}\ =\ t}}{Form 9: nested, const \{ p: \{ x \}, list: {[}first{]} \} = t}}\label{form-9-nested-const-p-x-list-first-t}

\textbf{Syntax:}
\passthrough{\lstinline!const \{ owner: \{ name \}, trades: [firstBuy] \} = position!}

\textbf{Example} (\passthrough{\lstinline!js/objects.mjs!}):

\begin{lstlisting}
const { owner: { name }, trades: [firstBuy] } = position
\end{lstlisting}

\textbf{Output:} \passthrough{\lstinline!form 9: Kamil 55!}

\textbf{Use it when:} the data I need is one or two levels deep.

\textbf{Watch out:} this does not create \passthrough{\lstinline!owner!}
or \passthrough{\lstinline!trades!} variables, they are only the path.
If \passthrough{\lstinline!owner!} is undefined it throws a TypeError.
Deep patterns get hard to read fast.

\subsubsection{\texorpdfstring{Form 10: computed key,
\texttt{const\ \{\ {[}key{]}:\ value\ \}\ =\ t}}{Form 10: computed key, const \{ {[}key{]}: value \} = t}}\label{form-10-computed-key-const-key-value-t}

\textbf{Syntax:}
\passthrough{\lstinline!const \{ [field]: chosen \} = position!}

\textbf{Example} (\passthrough{\lstinline!js/objects.mjs!}):

\begin{lstlisting}
const field = 'ticker'
const { [field]: chosen } = position
\end{lstlisting}

\textbf{Output:} \passthrough{\lstinline!form 10: ticker -> ASELS!}

\textbf{Use it when:} the property name is only known at run time (a
column chosen by the user).

\textbf{Watch out:} a computed key always needs a new name after the
colon (\passthrough{\lstinline!\{ [field] \}!} alone is a SyntaxError).
Plain version: \passthrough{\lstinline!const chosen = position[field]!}.

\subsubsection{\texorpdfstring{Form 11: patterns in parameters,
\texttt{function\ f(\{\ a\ =\ 1,\ b\ \}:\ Options\ =\ \{\})}}{Form 11: patterns in parameters, function f(\{ a = 1, b \}: Options = \{\})}}\label{form-11-patterns-in-parameters-function-f-a-1-b-options}

\textbf{Syntax:}
\passthrough{\lstinline!function placeOrder(\{ ticker, side = 'buy', qty = 1 \} = \{\})!}

\textbf{Example} (\passthrough{\lstinline!js/params.mjs!}):

\begin{lstlisting}
function placeOrder({ ticker, side = 'buy', qty = 1 } = {}) {
  return `${side} ${qty} x ${ticker}`
}
placeOrder({ ticker: 'BIMAS', qty: 5 })
placeOrder()
const spread = ([bid, ask]) => +(ask - bid).toFixed(2)
\end{lstlisting}

\textbf{Output:} \passthrough{\lstinline!form 11: buy 5 x BIMAS!},
\passthrough{\lstinline!form 11b: buy 1 x undefined!},
\passthrough{\lstinline!form 11c: 0.25!}

\textbf{Use it when:} a function has many options; the caller can pass
them in any order. React components are exactly this (P1).

\textbf{Watch out:} without \passthrough{\lstinline!= \{\}!} at the end,
\passthrough{\lstinline!placeOrder()!} throws ``Cannot destructure
property `ticker' of `undefined'\,''. In TypeScript the type goes after
the pattern: \passthrough{\lstinline!(\{ ... \}: OrderInput = \{\})!}.

\subsubsection{\texorpdfstring{Form 12: patterns in loops,
\texttt{for\ (const\ {[}k,\ v{]}\ of\ Object.entries(o))}}{Form 12: patterns in loops, for (const {[}k, v{]} of Object.entries(o))}}\label{form-12-patterns-in-loops-for-const-k-v-of-object.entrieso}

\textbf{Syntax:}
\passthrough{\lstinline!for (const [stock, weight] of weights)!}

\textbf{Example} (\passthrough{\lstinline!js/params.mjs!}):

\begin{lstlisting}
for (const [course, grade] of Object.entries(grades)) { ... }
const weights = new Map([['ASELS', 0.4], ['THYAO', 0.35], ['BIMAS', 0.25]])
for (const [stock, weight] of weights) { ... }
\end{lstlisting}

\textbf{Output:}

\begin{lstlisting}
form 12: MIS 233 AA
form 12: INTT 321 BA
form 12b: ASELS 0.4
form 12b: THYAO 0.35
form 12b: BIMAS 0.25
\end{lstlisting}

\textbf{Use it when:} looping over key/value pairs.

\textbf{Watch out:} \passthrough{\lstinline!Object.entries!} gives
\passthrough{\lstinline![key, value]!} arrays, so the pattern is an
\emph{array} pattern, not \passthrough{\lstinline!\{ key, value \}!}. A
Map already iterates as \passthrough{\lstinline![key, value]!} pairs.

\subsection{Pitfalls}\label{pitfalls}

All from \passthrough{\lstinline!js/pitfalls.mjs!}.

\begin{longtable}[]{@{}
  >{\raggedright\arraybackslash}p{(\columnwidth - 6\tabcolsep) * \real{0.0789}}
  >{\raggedright\arraybackslash}p{(\columnwidth - 6\tabcolsep) * \real{0.1579}}
  >{\raggedright\arraybackslash}p{(\columnwidth - 6\tabcolsep) * \real{0.3421}}
  >{\raggedright\arraybackslash}p{(\columnwidth - 6\tabcolsep) * \real{0.4211}}@{}}
\toprule\noalign{}
\begin{minipage}[b]{\linewidth}\raggedright
\#
\end{minipage} & \begin{minipage}[b]{\linewidth}\raggedright
Code
\end{minipage} & \begin{minipage}[b]{\linewidth}\raggedright
Real output
\end{minipage} & \begin{minipage}[b]{\linewidth}\raggedright
Why it happens
\end{minipage} \\
\midrule\noalign{}
\endhead
\bottomrule\noalign{}
\endlastfoot
1 & \passthrough{\lstinline!const \{ a \} = undefined!} &
\passthrough{\lstinline!pitfall 1: TypeError - Cannot destructure property 'a' of 'undefined' as it is undefined.!}
& Destructuring reads properties of the value, and
\passthrough{\lstinline!undefined!}/\passthrough{\lstinline!null!} have
no properties. \\
2 &
\passthrough{\lstinline!const \{ note = 'no note' \} = \{ note: null \}!}
& \passthrough{\lstinline!pitfall 2: null \\| no note!} & Defaults only
replace \passthrough{\lstinline!undefined!};
\passthrough{\lstinline!null!} is a real value, so it is kept. \\
3 & \passthrough{\lstinline!let a; \{ a \} = \{ a: 1 \}!} &
\passthrough{\lstinline!pitfall 3b: SyntaxError - Unexpected token '='!}
& A line starting with \passthrough{\lstinline!\{!} is parsed as a
block, so \passthrough{\lstinline!=!} makes no sense.
\passthrough{\lstinline!(\{ price \} = obj)!} works
(\passthrough{\lstinline!pitfall 3: 44.1!}). \\
4 &
\passthrough{\lstinline!const \{ meta \} = order; meta.note = 'changed'!}
& \passthrough{\lstinline!pitfall 4: changed!},
\passthrough{\lstinline!pitfall 4b: true!} & The variable points to the
same nested object; rest (\passthrough{\lstinline!...copy!}) is a
shallow copy too. \\
5 & \passthrough{\lstinline!const \{ b, a \} = \{a:1,b:2\}!} vs
\passthrough{\lstinline!const [y, x] = [1, 2]!} &
\passthrough{\lstinline!pitfall 5: object a = 1 b = 2 \\| array x = 2 y = 1!}
& Objects match by name, arrays by position. \\
\end{longtable}

\subsection{Props and children}\label{props-and-children}

Code in \passthrough{\lstinline!js/props.jsx!} and
\passthrough{\lstinline!ts/props.tsx!}, rendered with
\passthrough{\lstinline!renderToStaticMarkup!}.

\subsubsection{P1: props is one object}\label{p1-props-is-one-object}

\begin{lstlisting}
function CardA(props) { return <section><h2>{props.title}</h2>{props.children}</section> }
function CardB({ title, children }) { return <section><h2>{title}</h2>{children}</section> }
\end{lstlisting}

\begin{lstlisting}
P1 A: <section><h2>Portfolio</h2>3 positions</section>
P1 B: <section><h2>Portfolio</h2>3 positions</section>
P1 props object: [ 'ticker', 'qty', 'children' ]
\end{lstlisting}

React calls the component with ONE argument.
\passthrough{\lstinline!CardB!} is form 5 (properties by name) used in a
parameter, which is form 11. The \passthrough{\lstinline!Spy!} component
logs the keys of that one object: the JSX attributes plus
\passthrough{\lstinline!children!} for the text between the tags.

\subsubsection{P2: children is only a
prop}\label{p2-children-is-only-a-prop}

\begin{lstlisting}
P2 tags   : <section><h2>Note</h2>between tags</section>
P2 attr   : <section><h2>Note</h2>as attribute</section>
P2 both   : <section><h2>Note</h2>tags</section>
\end{lstlisting}

\passthrough{\lstinline!<CardB>between tags</CardB>!} and
\passthrough{\lstinline!<CardB children="as attribute" />!} give the
same HTML, so \passthrough{\lstinline!children!} is a normal prop with a
special JSX syntax. When I give both, the content between the tags wins,
because JSX puts it in \passthrough{\lstinline!children!} after the
attributes.

\subsubsection{P3: what children can be}\label{p3-what-children-can-be}

\begin{lstlisting}
P3 string   : <div>hello</div>
P3 number   : <div>42</div>
P3 zero     : <div>0</div>
P3 element  : <div><b>bold</b></div>
P3 several  : <div><i>a</i><i>b</i></div>
P3 array    : <div><span>x</span><span>y</span></div>
P3 undefined: <div></div>
P3 null     : <div></div>
P3 false    : <div></div>
P3 true     : <div></div>
P3 function : <p>Price: 44.1 TL</p>
P3 fn raw   : <div></div>
\end{lstlisting}

Strings, numbers (also \passthrough{\lstinline!0!}!) and elements
appear. \passthrough{\lstinline!undefined!},
\passthrough{\lstinline!null!}, \passthrough{\lstinline!true!} and
\passthrough{\lstinline!false!} render nothing, which is why
\passthrough{\lstinline!\{cond \&\& <X/>\}!} works, but
\passthrough{\lstinline!\{count \&\& ...\}!} prints
\passthrough{\lstinline!0!} when count is 0. A function cannot be
rendered directly: React prints nothing and warns on stderr ``Functions
are not valid as a React child''. It works only if the component
\emph{calls} \passthrough{\lstinline!children(price)!} (render prop).

\subsubsection{P4: separating children from the other
props}\label{p4-separating-children-from-the-other-props}

\begin{lstlisting}
P4 rest     : <section id="risk" class="card"><h3>Risk</h3>low</section>
P4 spread first : <section class="fixed">x</section>
P4 spread last  : <section class="mine">x</section>
P4 key      : <span>undefined</span>
\end{lstlisting}

\passthrough{\lstinline!(\{ title, children, ...rest \})!} takes out
what the component uses and passes the rest
(\passthrough{\lstinline!id!}, \passthrough{\lstinline!className!}) to
the DOM element. JSX attributes are applied left to right, so the one
that comes later wins:
\passthrough{\lstinline!\{...rest\} className="fixed"!} forces
\passthrough{\lstinline!fixed!},
\passthrough{\lstinline!className="fixed" \{...rest\}!} lets the caller
override it. \passthrough{\lstinline!key!} never reaches the component;
React warns ``\passthrough{\lstinline!key!} is not a prop. Trying to
access it will result in \passthrough{\lstinline!undefined!} being
returned''.

\subsubsection{P5: props or children?}\label{p5-props-or-children}

\begin{lstlisting}
P5 props    : <article><header>Weekly report</header><main>Up 2%</main><footer>Source: BIST</footer></article>
P5 children : <article><header>Weekly report</header><main>Up 2%</main><footer>Source: BIST</footer></article>
\end{lstlisting}

Same HTML, two designs. With props (\passthrough{\lstinline!header!},
\passthrough{\lstinline!body!}, \passthrough{\lstinline!footer!}) the
component controls the structure and the caller cannot forget a part or
put them in the wrong order. With children the caller controls
everything, which is more flexible but the component cannot enforce a
layout. I would use named props when there are fixed slots, and children
for a wrapper (card, modal) where the content is free.

\subsubsection{P6: defaults}\label{p6-defaults}

\begin{lstlisting}
P6 missing  : <span>[n/a]</span>
P6 undefined: <span>[n/a]</span>
P6 null     : <span>[]</span>
P6 value    : <span>[BUY]</span>
\end{lstlisting}

\passthrough{\lstinline!function Badge(\{ label = 'n/a' \})!} is form 7
in a parameter. A missing prop and
\passthrough{\lstinline!label=\{undefined\}!} get the default.
\passthrough{\lstinline!label=\{null\}!} does not, and since null
renders nothing the badge shows empty brackets on the screen.

\subsubsection{P7: types}\label{p7-types}

\begin{lstlisting}
P7 PropsWithChildren: <div><em>Watchlist</em>KCHOL, SISE</div>
P7 ComponentProps: <button type="button" disabled="">Buy</button>
\end{lstlisting}

\begin{itemize}
\tightlist
\item
  \passthrough{\lstinline!children?: ReactNode!} by hand: allows
  strings, numbers, elements, arrays, null, undefined, booleans. It
  rejects a function:
  \passthrough{\lstinline!TS2322: Type '() => string' is not assignable to type 'ReactNode'.!}
\item
  \passthrough{\lstinline!PropsWithChildren<\{ title: string \}>!} = my
  props plus \passthrough{\lstinline!children?: ReactNode!}, less
  typing.
\item
  \passthrough{\lstinline!ComponentProps<'button'>!} = every prop a real
  \passthrough{\lstinline!<button>!} accepts, so
  \passthrough{\lstinline!disabled!} works and a typo is caught:
  \passthrough{\lstinline!Property 'colour' does not exist on type '... ButtonHTMLAttributes<HTMLButtonElement>'. Did you mean 'color'?!}
\item
  Missing required prop:
  \passthrough{\lstinline!TS2741: Property 'title' is missing in type '\{ children: string; \}' but required in type 'CardProps'.!}
\item
  Wrong type:
  \passthrough{\lstinline!TS2322: Type 'number' is not assignable to type 'string'.!}
  In plain JS the same call just renders
  (\passthrough{\lstinline!P7 (js, no types): <section><h2>123</h2>4</section>!}).
\end{itemize}

The rejected lines are kept as comments in
\passthrough{\lstinline!ts/props.tsx!}, because the only allowed
suppression is the one in T5.

\subsection{TypeScript topics (T1 to
T5)}\label{typescript-topics-t1-to-t5}

From \passthrough{\lstinline!ts/types.ts!}:

\begin{itemize}
\tightlist
\item
  \textbf{T1}
  \passthrough{\lstinline!const \{ ticker: string \} = stock!} makes a
  variable called \passthrough{\lstinline!string!}
  (\passthrough{\lstinline!T1 rename trap: THYAO!}). The type goes after
  the pattern:
  \passthrough{\lstinline!const \{ ticker, price \}: \{ ticker: string; price: number \} = stock!}.
\item
  \textbf{T2} A tuple return \passthrough{\lstinline![number, string]!}
  gives each variable its own type. A normal array literal is
  \passthrough{\lstinline!(string | number)[]!}; with
  \passthrough{\lstinline!as const!} it becomes
  \passthrough{\lstinline!readonly [44.1, "ASELS"]!}, so
  \passthrough{\lstinline!p2!} has type \passthrough{\lstinline!44.1!}
  and \passthrough{\lstinline!t2!} has type
  \passthrough{\lstinline!"ASELS"!}.
\item
  \textbf{T3} With \passthrough{\lstinline!noUncheckedIndexedAccess!},
  \passthrough{\lstinline!const [first] = watchlist!} is
  \passthrough{\lstinline!string | undefined!}, so
  \passthrough{\lstinline!first.toLowerCase()!} does not compile until I
  check for undefined or add a default.
\item
  \textbf{T4} \passthrough{\lstinline!limit?: number!} →
  \passthrough{\lstinline!const \{ limit \}!} is
  \passthrough{\lstinline!number | undefined!}
  (\passthrough{\lstinline!T4: undefined 0!}); with
  \passthrough{\lstinline!= 0!} it becomes
  \passthrough{\lstinline!number!}.
\item
  \textbf{T5} \passthrough{\lstinline!const \{ last \} = quote!} when
  \passthrough{\lstinline!last!} is not in the type:
  \passthrough{\lstinline!Property 'last' does not exist on type '\{ bid: number; ask: number; \}'.!}
  (kept with \passthrough{\lstinline!// @ts-expect-error!}). Fix: add
  \passthrough{\lstinline!last?: number!} to the type, or give a
  default.
\end{itemize}

\subsection{JavaScript vs TypeScript}\label{javascript-vs-typescript}

\begin{longtable}[]{@{}
  >{\raggedright\arraybackslash}p{(\columnwidth - 6\tabcolsep) * \real{0.1017}}
  >{\raggedright\arraybackslash}p{(\columnwidth - 6\tabcolsep) * \real{0.1864}}
  >{\raggedright\arraybackslash}p{(\columnwidth - 6\tabcolsep) * \real{0.4068}}
  >{\raggedright\arraybackslash}p{(\columnwidth - 6\tabcolsep) * \real{0.3051}}@{}}
\toprule\noalign{}
\begin{minipage}[b]{\linewidth}\raggedright
Form
\end{minipage} & \begin{minipage}[b]{\linewidth}\raggedright
Same in TS
\end{minipage} & \begin{minipage}[b]{\linewidth}\raggedright
What I must add in TS
\end{minipage} & \begin{minipage}[b]{\linewidth}\raggedright
Error when wrong
\end{minipage} \\
\midrule\noalign{}
\endhead
\bottomrule\noalign{}
\endlastfoot
1 \passthrough{\lstinline![x, y] = f()!} & the pattern & a tuple return
type \passthrough{\lstinline![number, number]!} & without it each item
is \passthrough{\lstinline!number \\| undefined!}
(noUncheckedIndexedAccess) \\
2 \passthrough{\lstinline![, b, ...rest]!} & the pattern & nothing,
\passthrough{\lstinline!rest!} is \passthrough{\lstinline!string[]!} &
\passthrough{\lstinline!b!} is
\passthrough{\lstinline!string \\| undefined!},
\passthrough{\lstinline!b.toUpperCase()!}: ``TS18048: `b' is possibly
`undefined'.'' \\
3 swap & the pattern & compatible types for both variables & swapping a
string and a number: ``TS2322: Type `number' is not assignable to type
`string'.'' \\
4 array defaults & the pattern & optional tuple type
\passthrough{\lstinline![string?, number?]!} &
\passthrough{\lstinline!const [a4] = []!} with no type: ``TS2493: Tuple
type `{[}{]}' of length `0' has no element at index `0'.'' \\
5 \passthrough{\lstinline!\{ a, b \}!} & the pattern & object type
(often inferred) & asking for a field that is not in the type: TS2339
``Property \ldots{} does not exist'' \\
6 rename & the pattern & the type after the pattern, not inside &
\passthrough{\lstinline!\{ a: string \}!} silently renames to a variable
\passthrough{\lstinline!string!} (T1) \\
7 object defaults & the pattern & the field must be optional
(\passthrough{\lstinline!?!}) & the default removes
\passthrough{\lstinline!undefined!} from the type (T4) \\
8 rest & the pattern & nothing, rest type is computed & using a removed
field on the rest object: TS2339 \\
9 nested & the pattern & nested types & inner array item is
\passthrough{\lstinline!T \\| undefined!} \\
10 computed key & the pattern & key typed as
\passthrough{\lstinline!keyof T!} or a literal & with a plain
\passthrough{\lstinline!string!} key: ``TS2537: Type `\{ ticker: string;
\}' has no matching index signature for type `string'.'' \\
11 parameters & the pattern &
\passthrough{\lstinline!(\{ ... \}: Type = \{\})!} & missing required
option: TS2345 / TS2741 \\
12 loops & the pattern & \passthrough{\lstinline!Map<string, number>!}
or a \passthrough{\lstinline!Record!} &
\passthrough{\lstinline!Object.entries!} values are typed from the
record type \\
\end{longtable}

\subsection{Bonus 1: LaTeX}\label{bonus-1-latex}

The whole guide is also written in LaTeX:
\passthrough{\lstinline!latex/guide.tex!}, compiled to
\passthrough{\lstinline!latex/guide.pdf!}. It has a linked table of
contents, code typeset with the \passthrough{\lstinline!listings!}
package and real tables (\passthrough{\lstinline!booktabs!},
\passthrough{\lstinline!longtable!}). Commands I used (run twice so the
contents page is filled in):

\begin{lstlisting}
cd latex
pdflatex -interaction=nonstopmode guide.tex
pdflatex -interaction=nonstopmode guide.tex
\end{lstlisting}

\passthrough{\lstinline!latex/build.sh!} regenerates
\passthrough{\lstinline!guide.tex!} from this README with pandoc and
then runs the two commands above.

\subsection{Bonus 2: Depth}\label{bonus-2-depth}

Code in \passthrough{\lstinline!js/depth.mjs!} and
\passthrough{\lstinline!ts/depth.tsx!} (both run as part of
\passthrough{\lstinline!npm run all!}).

\subsubsection{D1: array patterns take any
iterable}\label{d1-array-patterns-take-any-iterable}

\begin{lstlisting}
D1 string: B I [ 'S', 'T' ]
D1 Set: ASELS
D1 Map: USDTRY 34.2
D1 generator: 40 41
\end{lstlisting}

A string gives its characters, a Set gives its unique values in
insertion order (so the duplicate \passthrough{\lstinline!ASELS!} is
only counted once), a Map gives \passthrough{\lstinline![key, value]!}
pairs (so I need a nested pattern \passthrough{\lstinline![[k, v]]!}),
and a generator gives what it yields.

\subsubsection{D2: what an array pattern really
calls}\label{d2-what-an-array-pattern-really-calls}

\begin{lstlisting}
D2 Symbol.iterator called
D2 next() call 1
D2 next() call 2
D2 return() called (iterator closed early)
D2 result: 10 20
D2 array-like: TypeError - {(intermediate value)(intermediate value)} is not iterable
\end{lstlisting}

The pattern calls \passthrough{\lstinline!value[Symbol.iterator]()!}
once, then \passthrough{\lstinline!next()!} only as many times as it has
positions (two here, not five). Because it stops early it calls
\passthrough{\lstinline!return()!} to close the iterator. An object with
\passthrough{\lstinline!0!} and \passthrough{\lstinline!length!} looks
like an array but has no \passthrough{\lstinline!Symbol.iterator!}, so
it throws ``is not iterable''. Array patterns are about the iteration
protocol, not about arrays.

\subsubsection{D3: order of evaluation}\label{d3-order-of-evaluation}

\begin{lstlisting}
D3 default computed for ask
D3 result: 39.9 0
D3 uses earlier variable: 1 2
D3 later variable: ReferenceError - Cannot access 'y' before initialization
\end{lstlisting}

Defaults are lazy: \passthrough{\lstinline!fallback('bid')!} never runs
because \passthrough{\lstinline!bid!} exists; only the missing
\passthrough{\lstinline!ask!} calls it. Defaults are evaluated left to
right, so \passthrough{\lstinline!high = low * 2!} can use
\passthrough{\lstinline!low!}, but
\passthrough{\lstinline![x = y, y = 1]!} throws a ReferenceError because
\passthrough{\lstinline!y!} is still in its temporal dead zone.

\subsubsection{D4: destructuring in other
places}\label{d4-destructuring-in-other-places}

\begin{lstlisting}
D4 catch: SyntaxError - Expected property name or '}' in JSON at position 2 (line 1 column 3)
D4 for...in gives keys: ASELS
D4 for...in gives keys: THYAO
D4 for...of entries: ASELS 120
D4 for...of entries: THYAO 40
D4 assignment targets: { min: 39.8, max: 43.5 }
D4 into array slots: [ 'x', 'y' ]
\end{lstlisting}

\passthrough{\lstinline!catch (\{ name, message \})!} takes the error
object apart directly. \passthrough{\lstinline!for...in!} gives only
keys (strings), \passthrough{\lstinline!for...of!} over
\passthrough{\lstinline!Object.entries!} gives pairs I can destructure.
In assignment (not declaration) the targets can be any property, like
\passthrough{\lstinline!target.min!} or
\passthrough{\lstinline!list[0]!}.

\subsubsection{D5: when destructuring makes code
worse}\label{d5-when-destructuring-makes-code-worse}

\begin{lstlisting}
D5 deep: BIST 22 11
D5 two steps: BIST 22 11
\end{lstlisting}

The one-line deep pattern
\passthrough{\lstinline!\{ meta: \{ source: \{ name: srcName, updated: \{ day, month \} \} \} \}!}
works, but it is hard to read, the names \passthrough{\lstinline!meta!},
\passthrough{\lstinline!source!}, \passthrough{\lstinline!updated!} are
only a path and are not variables, and if any level is missing it
throws. The two-step version is longer but each line is clear and easier
to debug. I would keep patterns to one level plus maybe one rename.

\subsubsection{D6 to D8: the same idea at the type
level}\label{d6-to-d8-the-same-idea-at-the-type-level}

\begin{lstlisting}
D6 First/Rest: ASELS [ 58.4, 1200 ]
D6 OwnerOf: { name: 'Kamil', account: 'Midas' }
D7 Parameters: THYAO 3 -> buy 3 x THYAO
D8 generic component: <ul><li>ASELS: 120</li><li>THYAO: 40</li></ul>
\end{lstlisting}

\begin{itemize}
\tightlist
\item
  \textbf{D6} \passthrough{\lstinline!First<T>!} and
  \passthrough{\lstinline!Rest<T>!} use \passthrough{\lstinline!infer!}
  inside a tuple pattern
  \passthrough{\lstinline![infer F, ...infer R]!}, which is form 2 for
  types. \passthrough{\lstinline!OwnerOf<T>!} does the same for an
  object type.
\item
  \textbf{D7} \passthrough{\lstinline!Parameters<typeof placeOrder>!}
  gives
  \passthrough{\lstinline![ticker: string, qty: number, side?: 'buy' | 'sell']!},
  so the parameters themselves are a tuple I can destructure and spread
  back in.
\item
  \textbf{D8} In the generic \passthrough{\lstinline!List<T>!},
  TypeScript infers \passthrough{\lstinline!T!} from
  \passthrough{\lstinline!items!}, so inside
  \passthrough{\lstinline!render=\{(\{ ticker, qty \}) => ...\}!} the
  pattern is fully typed without any annotation.
\end{itemize}

\subsection{Bonus 3: Real code}\label{bonus-3-real-code}

Five real cases (plus one more), with the file and line.

\subsubsection{\texorpdfstring{Case 1: A3 starter, \texttt{src/App.tsx}
line
37}{Case 1: A3 starter, src/App.tsx line 37}}\label{case-1-a3-starter-srcapp.tsx-line-37}

\textbf{Form:} 1

\begin{lstlisting}
const [tasks, dispatch] = useReducer(taskReducer, [])
\end{lstlisting}

Without destructuring:

\begin{lstlisting}
const r = useReducer(...); const tasks = r[0]; const dispatch = r[1]
\end{lstlisting}

\subsubsection{\texorpdfstring{Case 2: A3 starter,
\texttt{src/TaskRow.tsx} line
5}{Case 2: A3 starter, src/TaskRow.tsx line 5}}\label{case-2-a3-starter-srctaskrow.tsx-line-5}

\textbf{Form:} 11 (with 5)

\begin{lstlisting}
function TaskRow({ task, onNext, onDelete }: TaskRowProps)
\end{lstlisting}

Without destructuring:

\begin{lstlisting}
function TaskRow(props: TaskRowProps) { const task = props.task; ... }
\end{lstlisting}

\subsubsection{\texorpdfstring{Case 3: vite 8.3.8,
\texttt{node\_modules/vite/dist/node/chunks/node.js} line
6198}{Case 3: vite 8.3.8, node\_modules/vite/dist/node/chunks/node.js line 6198}}\label{case-3-vite-8.3.8-node_modulesvitedistnodechunksnode.js-line-6198}

\textbf{Form:} 11 + 7

\begin{lstlisting}
function resolveEnvPrefix({ envPrefix = "VITE_" })
\end{lstlisting}

Without destructuring:

\begin{lstlisting}
function resolveEnvPrefix(config) { let envPrefix = config.envPrefix === undefined ? "VITE_" : config.envPrefix }
\end{lstlisting}

\subsubsection{Case 4: vite 8.3.8, same file, line
1991}\label{case-4-vite-8.3.8-same-file-line-1991}

\textbf{Form:} 12 (in a normal \passthrough{\lstinline!for!})

\begin{lstlisting}
const entries = Object.entries(obj); // line 1988
const [key, value] = entries[i];    // line 1991, inside for (let i = 0; ...)
\end{lstlisting}

Without destructuring:

\begin{lstlisting}
const key = entries[i][0]; const value = entries[i][1]
\end{lstlisting}

\subsubsection{Case 5: vite 8.3.8, same file, line
22958}\label{case-5-vite-8.3.8-same-file-line-22958}

\textbf{Form:} 9

\begin{lstlisting}
const { EventTarget: { addEventListener, removeEventListener } } = require_event_target();
\end{lstlisting}

Without destructuring:

\begin{lstlisting}
const et = require_event_target().EventTarget; const addEventListener = et.addEventListener; ...
\end{lstlisting}

\subsubsection{\texorpdfstring{Case 6: shadcn/ui,
\texttt{apps/v4/registry/new-york-v4/ui/button.tsx} lines 40-49 (commit
\texttt{bfa64ff})}{Case 6: shadcn/ui, apps/v4/registry/new-york-v4/ui/button.tsx lines 40-49 (commit bfa64ff)}}\label{case-6-shadcnui-appsv4registrynew-york-v4uibutton.tsx-lines-40-49-commit-bfa64ff}

\textbf{Form:} 11 + 7 + 8

\begin{lstlisting}
function Button({ className, variant = "default", size = "default", asChild = false, ...props }: React.ComponentProps<"button"> & ...)
\end{lstlisting}

Without destructuring:

\begin{lstlisting}
function Button(allProps) { const variant = allProps.variant ?? "default"; ... const props = { ...allProps }; delete props.className; ... }
\end{lstlisting}

\textbf{Where it makes the code worse (case 5).} The nested pattern
hides the fact that \passthrough{\lstinline!EventTarget!} is an object;
a reader has to unpack two levels in their head, and if
\passthrough{\lstinline!EventTarget!} were missing the error would be a
confusing ``Cannot destructure property''. Two lines
(\passthrough{\lstinline!const \{ EventTarget \} = ...; const \{ addEventListener \} = EventTarget!})
say the same thing more clearly. Case 3 also has a small trap:
\passthrough{\lstinline!envPrefix!} is reassigned right after
(\passthrough{\lstinline!envPrefix = arraify(envPrefix)!}), so the same
name means two different types in one function.

\textbf{React props (case 6).} \passthrough{\lstinline!className!},
\passthrough{\lstinline!variant!}, \passthrough{\lstinline!size!} and
\passthrough{\lstinline!asChild!} are props the component uses itself;
\passthrough{\lstinline!...props!} collects every other prop of a real
\passthrough{\lstinline!<button>!} (\passthrough{\lstinline!onClick!},
\passthrough{\lstinline!disabled!}, \passthrough{\lstinline!type!},
\ldots) and also \passthrough{\lstinline!children!}, because children is
just one more prop. The button text between the tags is never named
here: it goes through \passthrough{\lstinline!\{...props\}!} onto
\passthrough{\lstinline!<Comp>!}, the same as P4 in my guide. Because
\passthrough{\lstinline!\{...props\}!} comes after
\passthrough{\lstinline!data-slot!} and
\passthrough{\lstinline!data-variant!}, a caller could override those,
but \passthrough{\lstinline!className!} is taken out first and merged
with \passthrough{\lstinline!cn(...)!}, so the caller's class is added,
not replacing the base classes.

\subsection{Bonus 4: Git repository}\label{bonus-4-git-repository}

https://github.com/kpostaagasi/mis233-assignment4 (public). The
repository has a \passthrough{\lstinline!.gitignore!} that leaves out
\passthrough{\lstinline!node\_modules!} and the LaTeX build files.

\subsection{My own log}\label{my-own-log}

TODO (Kamil): at least five real things that went wrong or surprised me
while I ran this project, with the exact error message and what I
learned. To be written by me after running the experiments myself.

\end{document}
